\documentclass[12pt,a4paper]{report}

\usepackage[utf8]{inputenc}
\usepackage[T1]{fontenc}
\usepackage[french]{babel}
\usepackage{geometry}
\usepackage{xcolor}
\usepackage{titlesec}
\usepackage{fancyhdr}
\usepackage{hyperref}
\usepackage{enumitem}
\usepackage{booktabs}
\usepackage{array}
\usepackage{mdframed}
\usepackage{setspace}
\usepackage{parskip}
\usepackage{colortbl}
\usepackage{tabularx}
\usepackage{listings}
\usepackage{float}
\usepackage{caption}
\usepackage{tikz}
\usetikzlibrary{arrows.meta, positioning, shapes.geometric}

\geometry{top=2.5cm, bottom=2.5cm, left=3cm, right=2.5cm}
\setlength{\headheight}{15pt}

% Colors
\definecolor{primaryblue}{RGB}{0, 71, 133}
\definecolor{accentblue}{RGB}{0, 120, 210}
\definecolor{lightgray}{RGB}{245, 245, 245}
\definecolor{darkgray}{RGB}{80, 80, 80}
\definecolor{alertred}{RGB}{180, 30, 30}
\definecolor{successgreen}{RGB}{30, 130, 70}
\definecolor{warnorange}{RGB}{200, 100, 0}
\definecolor{rowblue}{RGB}{220, 235, 250}
\definecolor{rowgreen}{RGB}{220, 245, 225}
\definecolor{rowred}{RGB}{250, 220, 220}
\definecolor{codebg}{RGB}{248, 248, 252}
\definecolor{codeframe}{RGB}{180, 200, 230}
\definecolor{codecomment}{RGB}{100, 130, 100}
\definecolor{codekeyword}{RGB}{0, 80, 160}
\definecolor{codestring}{RGB}{160, 60, 0}
\definecolor{codenumber}{RGB}{130, 130, 130}

% Section formatting
\titleformat{\chapter}[block]
  {\color{primaryblue}\Large\bfseries}{\thechapter.}{0.8em}{}
  [\color{primaryblue}\titlerule]
\titleformat{\section}{\color{accentblue}\large\bfseries}{\thesection}{0.6em}{}
\titleformat{\subsection}{\color{darkgray}\normalsize\bfseries}{\thesubsection}{0.5em}{}
\titlespacing*{\chapter}{0pt}{0pt}{1em}

% Header/Footer
\pagestyle{fancy}
\fancyhf{}
\fancyhead[L]{\small\color{darkgray}Rapport de Stage -- Semaines 11 \`{a} 15}
\fancyhead[R]{\small\color{darkgray}AROBASE SARL}
\fancyfoot[C]{\thepage}
\renewcommand{\headrulewidth}{0.4pt}

\hypersetup{colorlinks=true, linkcolor=primaryblue, urlcolor=accentblue,
  pdftitle={Rapport de Stage S11-S15 - AroBot RAG - Arobase SARL}}

% Info boxes
\newmdenv[backgroundcolor=lightgray, linecolor=accentblue, linewidth=2pt,
  leftline=true, topline=false, rightline=false, bottomline=false,
  innerleftmargin=12pt, innerrightmargin=8pt,
  innertopmargin=8pt, innerbottommargin=8pt]{infobox}

\newmdenv[backgroundcolor=alertred!8, linecolor=alertred, linewidth=2pt,
  leftline=true, topline=false, rightline=false, bottomline=false,
  innerleftmargin=12pt, innerrightmargin=8pt,
  innertopmargin=8pt, innerbottommargin=8pt]{warningbox}

\newmdenv[backgroundcolor=successgreen!8, linecolor=successgreen, linewidth=2pt,
  leftline=true, topline=false, rightline=false, bottomline=false,
  innerleftmargin=12pt, innerrightmargin=8pt,
  innertopmargin=8pt, innerbottommargin=8pt]{successbox}

\newmdenv[backgroundcolor=warnorange!10, linecolor=warnorange, linewidth=2pt,
  leftline=true, topline=false, rightline=false, bottomline=false,
  innerleftmargin=12pt, innerrightmargin=8pt,
  innertopmargin=8pt, innerbottommargin=8pt]{notebox}

% Python listings style
\lstdefinelanguage{Python}{
  keywords={import, from, def, class, return, if, elif, else, for, while,
             try, except, finally, with, as, in, not, and, or, is, lambda,
             True, False, None, self, super, async, await, yield, pass,
             break, continue, raise, global, nonlocal, del, assert},
  keywordstyle=\color{codekeyword}\bfseries,
  comment=[l]{\#},
  commentstyle=\color{codecomment}\itshape,
  stringstyle=\color{codestring},
  morestring=[b]",
  morestring=[b]',
  sensitive=true,
}

% JavaScript listings style
\lstdefinelanguage{JavaScript}{
  keywords={import, export, default, from, const, let, var, function, return,
             if, else, async, await, try, catch, new, null, true, false,
             process, require, module, this},
  keywordstyle=\color{codekeyword}\bfseries,
  comment=[l]{//},
  morecomment=[s]{/*}{*/},
  commentstyle=\color{codecomment}\itshape,
  stringstyle=\color{codestring},
  morestring=[b]",
  morestring=[b]',
  morestring=[b]`,
  sensitive=true,
}

\lstset{
  basicstyle=\ttfamily\footnotesize,
  backgroundcolor=\color{codebg},
  frame=single,
  framerule=0.5pt,
  rulecolor=\color{codeframe},
  numbers=left,
  numberstyle=\tiny\color{codenumber},
  numbersep=8pt,
  showstringspaces=false,
  breaklines=true,
  breakatwhitespace=true,
  tabsize=2,
  xleftmargin=15pt,
  xrightmargin=5pt,
  captionpos=b,
  aboveskip=10pt,
  belowskip=6pt,
}

% ============================================================
\begin{document}

% ==================== COVER PAGE ====================
\begin{titlepage}
  \centering\vspace*{1cm}
  {\color{primaryblue}\rule{\linewidth}{3pt}}\vspace{0.5cm}

  {\huge\bfseries\color{primaryblue} Rapport de Stage}\\[0.4cm]
  {\Large\color{darkgray} Semaines 11 \`{a} 15 -- Phase de D\'{e}veloppement RAG}\\[0.2cm]
  {\normalsize\color{darkgray} Justifications techniques, conception et pr\'{e}sentation du code produit}

  \vspace{0.5cm}{\color{primaryblue}\rule{\linewidth}{1.5pt}}\vspace{2cm}

  \begin{minipage}{0.48\textwidth}
    \centering\textbf{\large Entreprise d'accueil}\\[0.3cm]
    {\color{primaryblue}\large\bfseries AROBASE SARL}\\[0.2cm]
    \small Agence de cr\'{e}ation multim\'{e}dia\\
    \textit{``La solution multim\'{e}dia pour votre Communication !''}\\[0.3cm]
    19 rue cours Paoli\\
    20250 Corte -- Corse
  \end{minipage}
  \hfill
  \begin{minipage}{0.48\textwidth}
    \centering\textbf{\large Informations stage}\\[0.3cm]
    \begin{tabular}{ll}
      \small D\'{e}but du stage   & \small : F\'{e}vrier 2026 \\
      \small P\'{e}riode          & \small : S11 -- S15 \\
      \small Du                   & \small : 29 avril 2026 \\
      \small Au                   & \small : 5 juin 2026 \\
    \end{tabular}
  \end{minipage}

  \vspace{3cm}
  \fbox{\begin{minipage}{0.85\textwidth}
    \centering\vspace{0.3cm}
    {\large\bfseries Sujet}\\[0.3cm]
    \textit{D\'{e}veloppement d'un assistant conversationnel bas\'{e} sur RAG}\\
    \textit{pour le tourisme en Centre-Corse -- AroBot}\vspace{0.3cm}
  \end{minipage}}

  \vfill{\color{darkgray}\small Universit\'{e} de Corse Pasquale Paoli -- Master 2 Informatique -- Ann\'{e}e 2025-2026}
\end{titlepage}

\tableofcontents\newpage

% ============================================================
\chapter{Introduction -- Le projet AroBot RAG}

\section{Contexte du stage}

Dans le cadre du stage de Master 2, le projet confi\'{e} par AROBASE SARL consiste \`{a}
concevoir et d\'{e}velopper un \textbf{assistant conversationnel intelligent} d\'{e}di\'{e} au tourisme
en Centre-Corse. Cet assistant, baptis\'{e} \textbf{AroBot}, doit permettre aux visiteurs
d'obtenir des informations sur les h\'{e}bergements, activit\'{e}s et le patrimoine de la r\'{e}gion,
en interagissant en langage naturel.

\section{Description du projet}

AroBot repose sur l'architecture \textbf{RAG} (\textit{Retrieval-Augmented Generation}),
qui combine la recherche s\'{e}mantique dans une base documentaire locale avec la g\'{e}n\'{e}ration
de r\'{e}ponses par un mod\`{e}le de langage. Le syst\`{e}me est \textbf{enti\`{e}rement local} :
aucune donn\'{e}e n'est envoy\'{e}e vers des services cloud, ce qui garantit la confidentialit\'{e}
des documents des clients d'AROBASE.

\begin{center}
\renewcommand{\arraystretch}{1.3}
\begin{tabular}{>{\bfseries}p{4cm} p{9cm}}
  \toprule
  \rowcolor{rowblue}
  \textbf{Composant} & \textbf{Description} \\
  \midrule
  Corpus documentaire & Brochures PDF touristiques, guides d'h\'{e}bergement, sites web officiels \\
  Base vectorielle & ChromaDB (stockage local persistant des embeddings) \\
  Mod\`{e}le d'embedding & nomic-embed-text via Ollama (768 dimensions, multilingue) \\
  Mod\`{e}le de g\'{e}n\'{e}ration & Mistral via Ollama (r\'{e}ponses en langue naturelle) \\
  Interface & Widget JS int\'{e}grable + interface standalone Flask \\
  \bottomrule
\end{tabular}
\end{center}

\section{P\'{e}riode couverte -- Semaines 11 \`{a} 15}

\begin{center}
\renewcommand{\arraystretch}{1.4}
\begin{tabular}{>{\bfseries}p{1.5cm} p{3cm} p{8cm}}
  \toprule
  \rowcolor{rowblue}
  \textbf{Sem.} & \textbf{Dates} & \textbf{Th\`{e}me principal} \\
  \midrule
  S11 & 29 avr -- 5 mai & Analyse des besoins, choix des technologies, mise en place de l'environnement \\
  S12 & 6 -- 12 mai & D\'{e}veloppement des loaders PDF (7 modes) et strat\'{e}gies de chunking \\
  S13 & 13 -- 19 mai & Pipeline de retrieval hybride, scraping web, modes separate/mixed \\
  S14 & 20 -- 26 mai & Serveur Flask SSE, widget JS int\'{e}grable, interface standalone \\
  S15 & 27 mai -- 5 juin & Tests, \'{e}valuation Ragas, r\'{e}glage des param\`{e}tres RAG \\
  \bottomrule
\end{tabular}
\end{center}

% ============================================================
\chapter{Technologies utilis\'{e}es -- Justifications et comparatifs}

\section{\'{E}tape 1 -- Pr\'{e}sentation du projet et contraintes}

Avant de s\'{e}lectionner les technologies, il est n\'{e}cessaire d'identifier les contraintes
du projet. Ces contraintes permettront de d\'{e}gager les crit\`{e}res d'\'{e}valuation pertinents.

\begin{warningbox}
\textbf{Contraintes identifi\'{e}es :}
\begin{enumerate}[leftmargin=*]
  \item \textbf{D\'{e}ploiement local} : le syst\`{e}me doit fonctionner sans connexion internet
        (borne dans les offices de tourisme, zones \`{a} connectivit\'{e} limit\'{e}e en Corse).
  \item \textbf{Confidentialit\'{e}} : les documents touristiques des clients d'AROBASE ne
        doivent jamais transiter par des API cloud tierces.
  \item \textbf{Budget z\'{e}ro} : pas de coût r\'{e}current d'API (OpenAI, Pinecone...).
  \item \textbf{Corpus vari\'{e}} : brochures PDF \`{a} mise en page complexe (colonnes, tableaux
        de prix, photos) + pages web dynamiques.
  \item \textbf{Langue fran\c{c}aise} : corpus et requ\^{e}tes principalement en fran\c{c}ais,
        avec des terminologies corses et italiennes.
  \item \textbf{Int\'{e}gration l\'{e}g\`{e}re} : le widget doit pouvoir \^{e}tre int\'{e}gr\'{e} sur
        n'importe quel site client d'AROBASE via une simple balise \texttt{<script>}.
\end{enumerate}
\end{warningbox}

\section{\'{E}tape 2 -- Technologies candidates et crit\`{e}res d'\'{e}valuation}

Les contraintes font \'{e}merger trois \textbf{composants cl\'{e}s} \`{a} choisir :
la base de donn\'{e}es vectorielle, le mod\`{e}le de langage/embeddings,
et la biblioth\`{e}que d'extraction PDF. Pour chaque composant, entre 3 et 4 technologies
candidates ont \'{e}t\'{e} \'{e}valu\'{e}es.

Les \textbf{crit\`{e}res d'\'{e}valuation communs} (finalis\'{e}s apr\`{e}s analyse des technologies) sont :

\begin{center}
\renewcommand{\arraystretch}{1.3}
\begin{tabular}{>{\bfseries\small}p{0.5cm} >{\bfseries}p{4cm} p{8cm}}
  \toprule
  \rowcolor{rowblue}
  \textbf{N\textdegree} & \textbf{Crit\`{e}re} & \textbf{D\'{e}finition} \\
  \midrule
  C1 & Co\^{u}t de fonctionnement & Gratuit vs. abonnement/facturation \`{a} l'usage \\
  C2 & D\'{e}ploiement hors ligne & Fonctionnement sans connexion internet \\
  C3 & Confidentialit\'{e} des donn\'{e}es & Les donn\'{e}es restent sur la machine locale \\
  C4 & Qualit\'{e} de recherche & Pr\'{e}cision du retrieval (recall@k) \\
  C5 & Filtrage par m\'{e}tadonn\'{e}es & Possibilit\'{e} de filtrer par type, source, date... \\
  C6 & Int\'{e}gration Python & Facilit\'{e} d'utilisation depuis Python \\
  C7 & Support du fran\c{c}ais & Qualit\'{e} sur les textes en fran\c{c}ais \\
  C8 & Maintenabilit\'{e} & Communaut\'{e} active, documentation, licence open-source \\
  \bottomrule
\end{tabular}
\end{center}

\section{\'{E}tape 3 -- Comparatif : bases de donn\'{e}es vectorielles}

\textbf{Technologies candidates :} ChromaDB, Pinecone, Qdrant, Weaviate.

\begin{center}
\renewcommand{\arraystretch}{1.5}
\begin{tabularx}{\linewidth}{>{\bfseries\small}p{2.5cm} X X X X}
  \toprule
  \rowcolor{rowblue}
  \textbf{Crit\`{e}re} & \textbf{ChromaDB} & \textbf{Pinecone} & \textbf{Qdrant} & \textbf{Weaviate} \\
  \midrule
  C1 Co\^{u}t          & \cellcolor{rowgreen}\textbf{Gratuit}    & \cellcolor{rowred}\textbf{Payant}      & \cellcolor{rowgreen}\textbf{Gratuit}   & \cellcolor{warnorange!30}\textbf{Freemium} \\
  C2 Hors ligne         & \cellcolor{rowgreen}\textbf{Oui}        & \cellcolor{rowred}\textbf{Non}         & \cellcolor{rowgreen}\textbf{Oui}       & \cellcolor{warnorange!30}\textbf{Partiel}  \\
  C3 Confidentialit\'{e}& \cellcolor{rowgreen}\textbf{Totale}     & \cellcolor{rowred}\textbf{Cloud}       & \cellcolor{rowgreen}\textbf{Totale}    & \cellcolor{warnorange!30}\textbf{Option}   \\
  C4 Qualit\'{e}        & \cellcolor{warnorange!30}\textbf{Bonne} & \cellcolor{rowgreen}\textbf{Excellente}& \cellcolor{rowgreen}\textbf{Excellente}& \cellcolor{rowgreen}\textbf{Excellente}   \\
  C5 M\'{e}tadonn\'{e}es& \cellcolor{rowgreen}\textbf{Oui}       & \cellcolor{warnorange!30}\textbf{Limit\'{e}} & \cellcolor{rowgreen}\textbf{Oui}  & \cellcolor{rowgreen}\textbf{Oui}          \\
  C6 Python             & \cellcolor{rowgreen}\textbf{Native}     & \cellcolor{rowgreen}\textbf{SDK}       & \cellcolor{rowgreen}\textbf{SDK}       & \cellcolor{warnorange!30}\textbf{SDK}      \\
  C7 Fran\c{c}ais       & \cellcolor{rowgreen}\textbf{Oui}        & \cellcolor{rowgreen}\textbf{Oui}       & \cellcolor{rowgreen}\textbf{Oui}       & \cellcolor{rowgreen}\textbf{Oui}          \\
  C8 Maintenance        & \cellcolor{rowgreen}\textbf{Active}     & \cellcolor{rowgreen}\textbf{Active}    & \cellcolor{rowgreen}\textbf{Active}    & \cellcolor{rowgreen}\textbf{Active}       \\
  \bottomrule
\end{tabularx}
\end{center}

\begin{successbox}
\textbf{\'{E}tape 4 -- Choix retenu : ChromaDB.}
ChromaDB est la seule solution satisfaisant simultan\'{e}ment les contraintes C1, C2 et C3
(gratuit, hors ligne, local). Son int\'{e}gration Python est la plus simple du march\'{e}
(\texttt{chromadb.PersistentClient(path="./db")} en une ligne), et elle supporte nativement
le filtrage par m\'{e}tadonn\'{e}es avec une syntaxe \texttt{\$in}/\texttt{\$contains} directe.
Pinecone et Weaviate ont \'{e}t\'{e} \'{e}cart\'{e}s car ils n\'{e}cessitent une connexion cloud.
\end{successbox}

\section{\'{E}tape 3 -- Comparatif : mod\`{e}les de langage et d'embeddings}

\textbf{Technologies candidates :} Ollama (local), OpenAI API, Hugging Face Inference API.

\begin{center}
\renewcommand{\arraystretch}{1.5}
\begin{tabularx}{\linewidth}{>{\bfseries\small}p{2.5cm} X X X}
  \toprule
  \rowcolor{rowblue}
  \textbf{Crit\`{e}re} & \textbf{Ollama (local)} & \textbf{OpenAI API} & \textbf{Hugging Face API} \\
  \midrule
  C1 Co\^{u}t          & \cellcolor{rowgreen}\textbf{Gratuit}     & \cellcolor{rowred}\textbf{Pay\'{e}/token}  & \cellcolor{warnorange!30}\textbf{Freemium} \\
  C2 Hors ligne         & \cellcolor{rowgreen}\textbf{Oui}         & \cellcolor{rowred}\textbf{Non}             & \cellcolor{rowred}\textbf{Non}             \\
  C3 Confidentialit\'{e}& \cellcolor{rowgreen}\textbf{Totale}      & \cellcolor{rowred}\textbf{API externe}     & \cellcolor{rowred}\textbf{API externe}     \\
  C4 Qualit\'{e}        & \cellcolor{warnorange!30}\textbf{Bonne}  & \cellcolor{rowgreen}\textbf{Excellente}    & \cellcolor{warnorange!30}\textbf{Variable} \\
  C6 Python             & \cellcolor{rowgreen}\textbf{SDK natif}   & \cellcolor{rowgreen}\textbf{SDK}           & \cellcolor{rowgreen}\textbf{SDK}           \\
  C7 Fran\c{c}ais       & \cellcolor{rowgreen}\textbf{Mistral}     & \cellcolor{rowgreen}\textbf{GPT-4}         & \cellcolor{warnorange!30}\textbf{Variable} \\
  C8 Maintenance        & \cellcolor{rowgreen}\textbf{Active}      & \cellcolor{rowgreen}\textbf{Active}        & \cellcolor{rowgreen}\textbf{Active}        \\
  \bottomrule
\end{tabularx}
\end{center}

\begin{successbox}
\textbf{\'{E}tape 4 -- Choix retenu : Ollama + Mistral + nomic-embed-text.}
Ollama est la seule solution satisfaisant C1, C2 et C3 simultan\'{e}ment.
Le mod\`{e}le \textbf{Mistral} a \'{e}t\'{e} retenu pour la g\'{e}n\'{e}ration car il est optimis\'{e}
pour le fran\c{c}ais (d\'{e}velopp\'{e} par une entreprise fran\c{c}aise).
Le mod\`{e}le d'embedding \textbf{nomic-embed-text} (768 dimensions, multilingue) a \'{e}t\'{e}
pr\'{e}f\'{e}r\'{e} \`{a} \texttt{all-MiniLM-L6-v2} pour sa meilleure couverture du fran\c{c}ais
et sa capacit\'{e} de contexte plus large (8\,192 tokens).
\end{successbox}

\section{\'{E}tape 3 -- Comparatif : extraction de documents PDF}

Les brochures touristiques en PDF pr\'{e}sentent des d\'{e}fis sp\'{e}cifiques :
mises en page multi-colonnes, tableaux de prix, logos, photos.
Quatre strat\'{e}gies d'extraction ont \'{e}t\'{e} compar\'{e}es.

\begin{center}
\renewcommand{\arraystretch}{1.5}
\begin{tabularx}{\linewidth}{>{\bfseries\small}p{2.8cm} X X X X}
  \toprule
  \rowcolor{rowblue}
  \textbf{Crit\`{e}re} & \textbf{IBM Docling} & \textbf{PyMuPDF} & \textbf{PDFPlumber} & \textbf{Tesseract OCR} \\
  \midrule
  Layout complexe   & \cellcolor{rowgreen}\textbf{Excellent}     & \cellcolor{warnorange!30}\textbf{Moyen}     & \cellcolor{warnorange!30}\textbf{Bon}     & \cellcolor{rowred}\textbf{Faible}    \\
  Tableaux de prix  & \cellcolor{rowgreen}\textbf{Excellent}     & \cellcolor{rowred}\textbf{Faible}           & \cellcolor{rowgreen}\textbf{Excellent}   & \cellcolor{rowred}\textbf{Faible}    \\
  Documents scann\'{e}s& \cellcolor{rowred}\textbf{Non}           & \cellcolor{rowred}\textbf{Non}              & \cellcolor{rowred}\textbf{Non}           & \cellcolor{rowgreen}\textbf{Oui}     \\
  Sortie structur\'{e}e& \cellcolor{rowgreen}\textbf{Markdown}   & \cellcolor{warnorange!30}\textbf{Texte brut}& \cellcolor{warnorange!30}\textbf{Texte}  & \cellcolor{rowred}\textbf{Texte brut}\\
  Vitesse           & \cellcolor{rowred}\textbf{Lente (DL)}      & \cellcolor{rowgreen}\textbf{Rapide}         & \cellcolor{rowgreen}\textbf{Rapide}      & \cellcolor{rowred}\textbf{Tr\`{e}s lente} \\
  Open-source       & \cellcolor{rowgreen}\textbf{Oui}           & \cellcolor{rowgreen}\textbf{Oui}            & \cellcolor{rowgreen}\textbf{Oui}         & \cellcolor{rowgreen}\textbf{Oui}     \\
  \bottomrule
\end{tabularx}
\end{center}

\begin{successbox}
\textbf{\'{E}tape 4 -- Choix retenu : cascade Docling $\rightarrow$ PDFPlumber $\rightarrow$ PyMuPDF $\rightarrow$ Tesseract.}
Aucune technologie seule ne couvre tous les cas. La solution adopt\'{e}e est une
\textbf{cascade de fallbacks} configur\'{e}e par \texttt{PDF\_MODE} :
Docling (mode par d\'{e}faut pour les brochures touristiques complexes),
PDFPlumber (tableaux de prix tabulaires), PyMuPDF (documents simples, rapide),
Tesseract (pages scann\'{e}es uniquement). Docling est privil\'{e}gi\'{e} car il
utilise un mod\`{e}le de d\'{e}tection de layout par apprentissage profond
(IBM Research) et produit du Markdown structur\'{e} directement exploitable
par le MarkdownChunker.
\end{successbox}

\section{Veille technologique -- \'{E}volution des syst\`{e}mes RAG}

\subsection{Architecture RAG : de na\"{i}f \`{a} modulaire}

L'architecture RAG a consid\'{e}rablement \'{e}volu\'{e} depuis sa publication initiale
(Lewis et al., 2020). On distingue aujourd'hui trois g\'{e}n\'{e}rations :

\begin{center}
\renewcommand{\arraystretch}{1.4}
\begin{tabular}{>{\bfseries}p{3.5cm} p{3.5cm} p{6cm}}
  \toprule
  \rowcolor{rowblue}
  \textbf{G\'{e}n\'{e}ration} & \textbf{Approche} & \textbf{Innovation} \\
  \midrule
  Naive RAG (2020) & Chunk $\rightarrow$ Embed $\rightarrow$ Search $\rightarrow$ Generate & Concept de base \\
  Advanced RAG (2023) & Query rewriting, HyDE, Reranking & Am\'{e}lioration de la pr\'{e}cision \\
  Modular RAG (2024) & Graph RAG, Agentic RAG, Self-RAG & Flexibilit\'{e} et r\'{e}flexivit\'{e} \\
  \bottomrule
\end{tabular}
\end{center}

\subsection{Techniques impl\'{e}ment\'{e}es dans AroBot}

Le projet AroBot int\`{e}gre plusieurs techniques d'Advanced RAG :

\begin{itemize}[leftmargin=*]
  \item \textbf{Query Rewriting} : reformulation automatique des questions de suivi
        (\og{}et lui ?\fg{} $\rightarrow$ question autonome) via un appel LLM d\'{e}di\'{e}.
  \item \textbf{Hybrid Search} : combinaison de la recherche dense (embeddings)
        et de la recherche par mots-cl\'{e}s (entit\'{e}s nomm\'{e}es comme \og{}H\^{o}tel Sampiero\fg{}).
  \item \textbf{Parent-Document Retrieval} : recherche sur de petits chunks pr\'{e}cis,
        r\'{e}cup\'{e}ration de grands chunks riches pour le LLM.
  \item \textbf{Distance Filtering} : rejet des chunks trop \'{e}loign\'{e}s s\'{e}mantiquement
        (\texttt{MAX\_DISTANCE = 0.9}) pour \'{e}viter les hallucinations.
\end{itemize}

\subsection{Tendances 2025 surveill\'{e}es}

\begin{itemize}[leftmargin=*]
  \item \textbf{Llama 3.x / Phi-3 / Gemma 2} : alternatives \`{a} Mistral via Ollama,
        potentiellement plus performantes sur le fran\c{c}ais.
  \item \textbf{ColBERT / SPLADE} : mod\`{e}les de recherche hybride end-to-end plus pr\'{e}cis
        que la combinaison embedding + keyword.
  \item \textbf{Ragas 0.2} : \'{e}volution du framework d'\'{e}valuation utilis\'{e} dans ce projet
        (m\'{e}triques faithfulness, context\_precision, answer\_relevancy).
  \item \textbf{Multimodal RAG} : ChromaDB pr\'{e}pare le support des embeddings d'images,
        pertinent pour les brochures touristiques avec photos.
\end{itemize}

% ============================================================
\chapter{Conception -- Architecture et diagrammes}

\section{Architecture g\'{e}n\'{e}rale du syst\`{e}me}

Le syst\`{e}me est organis\'{e} en quatre couches ind\'{e}pendantes :

\begin{center}
\begin{tikzpicture}[
  box/.style={rectangle, draw=accentblue, fill=rowblue,
              minimum width=5cm, minimum height=0.75cm,
              text centered, font=\small\bfseries, rounded corners=2pt},
  sbox/.style={rectangle, draw=successgreen!80, fill=rowgreen,
               minimum width=2.3cm, minimum height=0.7cm,
               text centered, font=\small, rounded corners=2pt},
  arr/.style={->, >=Stealth, thick, color=darkgray},
  lbl/.style={font=\footnotesize\color{darkgray}}
]
  \node[box] (front)  at (0, 0)    {Couche Pr\'{e}sentation -- Widget.js / index.html};
  \node[box] (api)    at (0, -1.6) {Couche API -- Flask (server.py)};
  \node[box] (rag)    at (0, -3.2) {Couche RAG -- Pipeline (legal\_rag/)};
  \node[sbox](chroma) at (-2.5,-4.9){ChromaDB};
  \node[sbox](ollama) at ( 2.5,-4.9){Ollama};

  \draw[arr] (front)  -- node[lbl,right]{POST /ask (JSON + session\_id)} (api);
  \draw[arr] (api)    -- node[lbl,right]{search() + stream\_answer()} (rag);
  \draw[arr] (rag)    -- node[lbl,above left]{query embedding} (chroma);
  \draw[arr] (chroma) -- node[lbl,below left]{chunks + distances} (rag);
  \draw[arr] (rag)    -- node[lbl,above right]{prompt + contexte} (ollama);
  \draw[arr] (ollama) -- node[lbl,below right]{tokens SSE} (rag);
  \draw[arr] (rag)    to[bend left=25] node[lbl,left]{stream SSE} (api);
  \draw[arr] (api)    to[bend left=25] node[lbl,left]{tokens + sources} (front);
\end{tikzpicture}
\end{center}

\section{Diagramme de s\'{e}quence -- Traitement d'une requ\^{e}te}

Le tableau suivant d\'{e}crit le flux complet d'une question de l'utilisateur \`{a} la r\'{e}ponse finale :

\begin{center}
\renewcommand{\arraystretch}{1.4}
\begin{tabularx}{\linewidth}{>{\small\bfseries}p{0.4cm} p{2.2cm} >{\small\ttfamily}p{3.5cm} X}
  \toprule
  \rowcolor{rowblue}
  \textbf{\#} & \textbf{Acteur} & \textbf{Action} & \textbf{Description} \\
  \midrule
  1 & Utilisateur      & saisit question     & L'utilisateur tape sa question dans le widget \\
  2 & Widget.js        & POST /ask           & Envoi JSON \{query, session\_id, key\} au serveur Flask \\
  3 & Flask            & rewrite\_query()    & Reformulation si question de suivi (appel Ollama) \\
  4 & Flask            & pipeline.search()   & Recherche hybride dans ChromaDB (dense + keyword) \\
  5 & ChromaDB         & retourne chunks     & Top-12 chunks triés par distance L2 \\
  6 & Flask            & build\_context()    & Filtre les chunks (distance $\leq$ 0.9), formate le contexte \\
  7 & Flask            & stream\_answer()    & Construction du prompt RAG + appel Ollama en streaming \\
  8 & Ollama           & yield token         & Génère la réponse token par token (temperature=0.0) \\
  9 & Flask            & SSE yield           & Transmet chaque token via Server-Sent Events \\
  10 & Widget.js       & affiche token       & Mise à jour progressive de la bulle de réponse \\
  11 & Flask           & done + sources      & Fin du stream, envoi des sources citables \\
  12 & Widget.js       & affiche sources     & Puces cliquables sous la réponse \\
  \bottomrule
\end{tabularx}
\end{center}

\section{Strat\'{e}gies de d\'{e}coupage (Chunking)}

Le choix de la strat\'{e}gie de d\'{e}coupage est d\'{e}terminant pour la qualit\'{e} du retrieval.
Trois approches ont \'{e}t\'{e} impl\'{e}ment\'{e}es, chacune adapt\'{e}e \`{a} un type de contenu :

\begin{center}
\renewcommand{\arraystretch}{1.5}
\begin{tabularx}{\linewidth}{>{\bfseries}p{3.2cm} X X}
  \toprule
  \rowcolor{rowblue}
  \textbf{Strat\'{e}gie} & \textbf{StructuralChunker} & \textbf{MarkdownChunker} \\
  \midrule
  Principe de d\'{e}coupe & Paragraphes (\textbackslash n\textbackslash n) avec fusion intelligente & Titres Markdown (\#\#) issus de Docling \\
  Chevauchement & 400 caract\`{e}res (2 phrases) & 0 (fronti\`{e}res s\'{e}mantiques nettes) \\
  Taille max & 2\,000 caract\`{e}res & 2\,000 caract\`{e}res \\
  Taille min & 300 caract\`{e}res & 100 caract\`{e}res \\
  Corpus cible & PDF texte brut, XML, JSON & Brochures touristiques (Docling) \\
  Avantage & Robuste sur tout type de texte & Chaque h\^{o}tel/activit\'{e} = 1 chunk \\
  \bottomrule
\end{tabularx}
\end{center}

\vspace{0.5em}
\begin{notebox}
\textbf{Strat\'{e}gie Parent-Document Retriever :} une troisi\`{e}me strat\'{e}gie optionnelle
utilise deux collections ChromaDB -- des \og{}enfants\fg{} (800 chars, recherch\'{e}s par embedding)
et des \og{}parents\fg{} (2\,500 chars, envoy\'{e}s au LLM). Cette approche r\'{e}sout le
\textit{trilemme du chunking} : pr\'{e}cision de recherche ET richesse de contexte.
\end{notebox}

\section{Cahier des charges -- Exigences fonctionnelles et non-fonctionnelles}

\begin{center}
\renewcommand{\arraystretch}{1.4}
\begin{tabularx}{\linewidth}{>{\bfseries\small}p{0.6cm} >{\bfseries}p{4cm} X p{2.5cm}}
  \toprule
  \rowcolor{rowblue}
  \textbf{ID} & \textbf{Exigence} & \textbf{Description} & \textbf{Statut} \\
  \midrule
  \rowcolor{rowgreen!50}
  F1 & R\'{e}ponse en langage naturel & L'assistant r\'{e}pond en fran\c{c}ais \`{a} partir du corpus & \textcolor{successgreen}{\textbf{Impl\'{e}ment\'{e}}} \\
  \rowcolor{rowgreen!50}
  F2 & Sources cit\'{e}es & Chaque r\'{e}ponse indique les documents sources & \textcolor{successgreen}{\textbf{Impl\'{e}ment\'{e}}} \\
  \rowcolor{rowgreen!50}
  F3 & Streaming temps r\'{e}el & Affichage progressif token par token & \textcolor{successgreen}{\textbf{Impl\'{e}ment\'{e}}} \\
  \rowcolor{rowgreen!50}
  F4 & Widget int\'{e}grable & Script JS embarquable sur tout site & \textcolor{successgreen}{\textbf{Impl\'{e}ment\'{e}}} \\
  \rowcolor{rowgreen!50}
  F5 & Multi-client & Clés API distinctes, corpus s\'{e}par\'{e}s & \textcolor{successgreen}{\textbf{Impl\'{e}ment\'{e}}} \\
  \rowcolor{warnorange!20}
  F6 & Historique conversation & Contexte des 3 derniers \'{e}changes & \textcolor{warnorange}{\textbf{Partiel}} \\
  \midrule
  \rowcolor{rowgreen!50}
  NF1 & D\'{e}ploiement 100\% local & Z\'{e}ro appel cloud en production & \textcolor{successgreen}{\textbf{Impl\'{e}ment\'{e}}} \\
  \rowcolor{rowgreen!50}
  NF2 & Temps de r\'{e}ponse $<$ 3s & Premier token en moins de 3 secondes & \textcolor{successgreen}{\textbf{Impl\'{e}ment\'{e}}} \\
  \rowcolor{rowgreen!50}
  NF3 & Anti-hallucination & Refus si aucun passage pertinent & \textcolor{successgreen}{\textbf{Impl\'{e}ment\'{e}}} \\
  \rowcolor{rowred!50}
  NF4 & \'{E}valuation automatique & M\'{e}triques Ragas (faithfulness, precision) & \textcolor{alertred}{\textbf{En cours}} \\
  \bottomrule
\end{tabularx}
\end{center}

% ============================================================
\chapter{Code produit -- Exemples comment\'{e}s}

\section{Contexte g\'{e}n\'{e}ral}

Les quatre fichiers pr\'{e}sent\'{e}s dans ce chapitre constituent le c\oe{}ur technique du syst\`{e}me.
Ils illustrent les d\'{e}cisions de conception valid\'{e}es avec le ma\^{i}tre de stage.

\begin{center}
\renewcommand{\arraystretch}{1.4}
\begin{tabular}{p{0.5cm} >{\ttfamily}p{3.5cm} p{8.5cm}}
  \toprule
  \rowcolor{rowblue}
  \textbf{N\textdegree} & \textbf{Fichier} & \textbf{R\^{o}le} \\
  \midrule
  1 & legal\_rag/pipeline.py   & Orchestration de l'ingestion -- s\'{e}lection dynamique du chunker \\
  2 & legal\_rag/generation.py & Construction du prompt RAG et g\'{e}n\'{e}ration en streaming \\
  3 & server.py                & Route Flask \texttt{/ask} avec Server-Sent Events \\
  4 & static/widget.js         & Client JavaScript -- envoi de requ\^{e}te et lecture du stream \\
  \bottomrule
\end{tabular}
\end{center}

% ============================================================
\section{pipeline.py -- S\'{e}lection de la strat\'{e}gie de chunking}

\subsection{Contexte et r\^{o}le}

\texttt{IngestionPipeline} est le point d'entr\'{e}e de l'ingestion documentaire. Son constructeur
illustre la \textbf{strat\'{e}gie de s\'{e}lection dynamique} du chunker selon le mode d'extraction PDF
configur\'{e}, ainsi que le choix de la strat\'{e}gie d'indexation (recursive vs. parent-child).

\subsection{Code comment\'{e}}

\begin{lstlisting}[language=Python, caption={pipeline.py -- Constructeur d'IngestionPipeline}]
class IngestionPipeline:
    """Pipeline complet d'ingestion pour tout type de corpus."""

    def __init__(self, collection_name: str = "legal_corpus_m2_tp",
                 retriever_type: str = "recursive"):
        self.retriever_type = retriever_type
        self.collection_name = collection_name

        # Si PDF_MODE produit du Markdown structure (docling, pymupdf4llm),
        # on utilise MarkdownChunker qui decoupe sur les titres ##.
        # Chaque section "## Hotel Le Paradis" devient un chunk autonome.
        # Si PDF_MODE produit du texte brut, on utilise StructuralChunker
        # qui decoupe par paragraphes avec un chevauchement de 400 chars.
        if PDF_MODE in ("docling", "markdown"):
            self.chunker = MarkdownChunker(
                max_chunk_size=MAX_CHUNK_SIZE,  # 2 000 caracteres par chunk
                min_chunk_size=100,             # Evite les titres sans contenu
                overlap=0                        # Frontiere semantique nette sur ##
            )
        else:
            self.chunker = StructuralChunker(
                max_chunk_size=MAX_CHUNK_SIZE,
                min_chunk_size=300,
                overlap=400  # 400 chars = ~2 phrases pour preserver le contexte
            )

        # Choix de la strategie d'indexation/retrieval
        # "parent-child" : 2 collections (enfants pour recherche, parents pour LLM)
        # "recursive"    : 1 collection, retrieval direct (defaut)
        if self.retriever_type == "parent-child":
            self.parent_retriever = ParentDocumentRetriever(
                collection_name_children=f"{collection_name}_children",
                collection_name_parents=f"{collection_name}_parents"
            )
        else:
            # Mode recursif : une seule collection ChromaDB persistante
            self.indexer = CorpusIndexer(collection_name=collection_name)
\end{lstlisting}

\textbf{Explication :} La s\'{e}lection conditionnelle du chunker est un \textbf{pattern Strategy}
simplifi\'{e} : le comportement de d\'{e}coupage est encapsul\'{e} dans deux classes interchangeables.
Cette approche permet de changer de strat\'{e}gie d'extraction (et donc de chunking) en modifiant
une seule variable \texttt{PDF\_MODE} dans \texttt{config.py}, sans toucher au reste du pipeline.
Le chevauchement de 400 caract\`{e}res dans \texttt{StructuralChunker} correspond \`{a} environ
deux phrases, ce qui garantit que le contexte d'une phrase coup\'{e}e \`{a} la fronti\`{e}re d'un chunk
est pr\'{e}sent dans le chunk suivant.

% ============================================================
\section{generation.py -- Prompt RAG et streaming}

\subsection{Contexte et r\^{o}le}

\texttt{AnswerGenerator} encapsule toute la logique de g\'{e}n\'{e}ration. Deux m\'{e}thodes sont
pr\'{e}sent\'{e}es : \texttt{build\_prompt()} qui construit le prompt syst\`{e}me inject\'{e} avec le
contexte retrieval, et \texttt{stream\_answer()} qui appelle Ollama en mode streaming.

\subsection{Code comment\'{e} -- Construction du prompt}

\begin{lstlisting}[language=Python, caption={generation.py -- build\_prompt() : injection du contexte RAG}]
def build_prompt(self, query: str, results: dict) -> str:
    """Construit le prompt complet envoye au LLM."""

    # Construit le contexte a partir des chunks recuperes par ChromaDB.
    # _build_context() filtre les chunks trop eloignes (distance > MAX_DISTANCE)
    # et formate chaque passage avec sa source.
    context = self._build_context(results)

    # Si aucun passage pertinent n'a ete trouve, on refuse de repondre.
    # Cela evite que le LLM "invente" une reponse sans source.
    if not context:
        return ""

    # Construction du prompt final en trois parties :
    # 1. Instruction systeme (personnalite du domaine)
    # 2. Passages recuperes (contexte RAG)
    # 3. Question de l'utilisateur
    return f"""{self.system_intro}

Regles strictes :
- Reponds UNIQUEMENT a ce qui est demande, de facon courte et directe.
- Utilise SEULEMENT les passages qui repondent precisement a la question.
- Ne cite jamais les sources dans ta reponse, elles sont affichees separement.
- Si l'information est absente : "Je n'ai pas cette information dans les documents."
- Jamais de connaissance generale, jamais d'invention.

PASSAGES :
{context}

QUESTION : {query}
REPONSE (courte et directe) :"""
\end{lstlisting}

\subsection{Code comment\'{e} -- Streaming token par token}

\begin{lstlisting}[language=Python, caption={generation.py -- stream\_answer() : generation en streaming}]
def stream_answer(self, query: str, results: dict, history: list = None):
    """Genere la reponse token par token via Ollama (generateur Python)."""

    prompt = self.build_prompt(query, results)
    if not prompt:
        return  # Stoppe si aucun contexte pertinent

    # Construction du contexte de conversation (historique des 3 derniers echanges).
    # Le message systeme definit la personnalite de l'assistant.
    messages = [{"role": "system", "content": self.system_intro}]
    for msg in (history or []):
        messages.append({"role": msg["role"], "content": msg["content"]})
    messages.append({"role": "user", "content": prompt})

    # Appel Ollama en mode streaming (stream=True).
    # temperature=0.0 assure des reponses deterministes et reproductibles.
    # C'est critique pour un assistant factuel : on ne veut pas de variabilite.
    stream = ollama.chat(
        model=self.model,
        messages=messages,
        stream=True,
        options={"temperature": 0.0}
    )

    # Yield chaque token au fur et a mesure de la generation.
    # Flask transmettra ces tokens via Server-Sent Events au client.
    for chunk in stream:
        token = chunk.message.content
        if token:
            yield token
\end{lstlisting}

\textbf{Explication :} La fonction \texttt{stream\_answer()} est un \textbf{g\'{e}n\'{e}rateur Python}
(utilisant \texttt{yield}). Cela permet \`{a} Flask de transmettre chaque token au client
imm\'{e}diatement, sans attendre la fin de la g\'{e}n\'{e}ration. L'utilisateur voit ainsi la
r\'{e}ponse appara\^{i}tre progressivement, ce qui am\'{e}liore consid\'{e}rablement l'exp\'{e}rience
ressentie. La valeur \texttt{temperature=0.0} est fondamentale : elle rend les r\'{e}ponses
d\'{e}terministes (toujours la m\^{e}me r\'{e}ponse pour la m\^{e}me question et le m\^{e}me contexte),
ce qui est essentiel pour un assistant factuel bas\'{e} sur des documents touristiques.

% ============================================================
\section{server.py -- Route /ask et Server-Sent Events}

\subsection{Contexte et r\^{o}le}

La route \texttt{/ask} est le \textbf{point d'entr\'{e}e principal} de l'API. Elle orchestre
l'ensemble du pipeline RAG (reformulation, recherche, g\'{e}n\'{e}ration) et utilise le protocole
\textbf{Server-Sent Events (SSE)} pour transmettre la r\'{e}ponse au client en temps r\'{e}el.

\subsection{Code comment\'{e}}

\begin{lstlisting}[language=Python, caption={server.py -- Route /ask avec streaming SSE}]
@app.route("/ask", methods=["POST"])
def ask():
    data       = request.json or {}
    query      = data.get("query", "").strip()
    api_key    = data.get("key", "")
    session_id = data.get("session_id", "")

    if not query:
        return jsonify({"error": "Question vide"}), 400

    # Multi-client : chaque cle API pointe vers une collection ChromaDB distincte.
    # Un client sans cle utilise le pipeline local par defaut.
    if api_key and api_key not in CLIENTS:
        return jsonify({"error": "Cle API invalide"}), 403

    active_pipeline = PIPELINES.get(api_key, LOCAL_PIPELINE)

    # Recuperation des 6 derniers messages de la session (3 echanges max).
    history = get_history(session_id) if session_id else []

    # Reformulation intelligente des questions de suivi :
    # "et lui ?" -> "Quels sont les horaires de l'hotel Sampiero ?"
    search_query = generator.rewrite_query(query, history)
    results = active_pipeline.search(query=search_query, n_results=12)

    no_ids = not results or not results.get("ids") or not results["ids"][0]
    if no_ids:
        def no_res():
            yield f"data: {json.dumps({'error': 'no_results'})}\n\n"
        return Response(stream_with_context(no_res()), mimetype="text/event-stream")

    context = generator.build_context(results)
    if not context:
        def no_ctx():
            yield f"data: {json.dumps({'error': 'no_context'})}\n\n"
        return Response(stream_with_context(no_ctx()), mimetype="text/event-stream")

    # Preparation des sources citables (fichiers locaux ou URLs web).
    # Deduplication par nom de fichier, filtrage par MAX_DISTANCE.
    sources = []
    seen = set()
    distances = results.get("distances", [[]])[0]
    for i, meta in enumerate(results["metadatas"][0]):
        dist = distances[i] if i < len(distances) else None
        if dist is not None and dist > MAX_DISTANCE:
            continue  # Chunk trop eloigne semantiquement : on ne le cite pas
        raw = meta.get("source_file", "")
        source_type = meta.get("source_type", "")
        filename = raw.split("/")[-1].split("\\")[-1] or raw or "source"
        url = raw if source_type == "web" else f"/documents/local/{filename}"
        if filename in seen:
            continue
        seen.add(filename)
        sources.append({"label": filename, "url": url, "dist": round(dist, 2)})

    def generate():
        full = ""
        # Stream chaque token de la reponse vers le client
        for token in generator.stream_answer(query, results, history=history):
            full += token
            # Format SSE : "data: {json}\n\n"
            yield f"data: {json.dumps({'token': token})}\n\n"
        # Fin du stream : envoi des sources et sauvegarde en session
        show_sources = sources if "je n'ai pas cette information" not in full.lower() else []
        if session_id and show_sources:
            save_exchange(session_id, query, full)
        yield f"data: {json.dumps({'done': True, 'sources': show_sources})}\n\n"

    # mimetype="text/event-stream" signale au navigateur qu'il s'agit de SSE.
    # stream_with_context() preserves le contexte Flask pendant le streaming.
    return Response(stream_with_context(generate()), mimetype="text/event-stream")
\end{lstlisting}

\textbf{Explication :} Le protocole \textbf{Server-Sent Events (SSE)} est pr\'{e}f\'{e}r\'{e} aux
WebSockets pour ce cas d'usage car il est unidirectionnel (serveur vers client),
plus simple \`{a} impl\'{e}menter avec Flask, et nativement support\'{e} par les navigateurs
modernes sans biblioth\`{e}que suppl\'{e}mentaire. Chaque \'{e}v\'{e}nement SSE a le format
\texttt{data: \{json\}{\textbackslash}n{\textbackslash}n}. Deux types de messages sont envoy\'{e}s :
\texttt{\{"token": "..."\}} pendant la g\'{e}n\'{e}ration, et
\texttt{\{"done": true, "sources": [...]\}} \`{a} la fin. La condition sur
\texttt{"je n'ai pas cette information"} \'{e}vite d'afficher des sources quand l'assistant
reconn\^{a}it lui-m\^{e}me qu'il ne peut pas r\'{e}pondre.

% ============================================================
\section{widget.js -- Envoi de requ\^{e}te et lecture du stream}

\subsection{Contexte et r\^{o}le}

\texttt{widget.js} est un script JavaScript auto-ex\'{e}cutant (\textit{IIFE -- Immediately
Invoked Function Expression}) inject\'{e} sur n'importe quelle page cliente via une balise
\texttt{<script>}. La fonction \texttt{sendQuery()} g\`{e}re l'envoi de la requ\^{e}te,
la lecture du stream SSE, et l'affichage progressif de la r\'{e}ponse.

\subsection{Code comment\'{e}}

\begin{lstlisting}[language=JavaScript, caption={widget.js -- sendQuery() avec AbortController et lecture SSE}]
async function sendQuery() {
  // Si une generation est deja en cours, on l'interrompt immediatement.
  // L'AbortController permet d'annuler le fetch cote client.
  // Cela evite les doublons si l'utilisateur clique trop vite.
  if (abortCtrl) { abortCtrl.abort(); return; }

  const q = inputEl.value.trim();
  if (!q) return;

  addBubble("user").textContent = q;  // Bulle utilisateur, affichage immediat
  inputEl.value = "";
  setStopMode(true);       // Remplace le bouton Envoyer par un carre Stop
  const typing = addTyping(); // Trois points animes pendant l'attente

  abortCtrl = new AbortController(); // Cree un controleur d'annulation
  try {
    const res = await fetch(`${HOST}/ask`, {
      method: "POST",
      headers: { "Content-Type": "application/json" },
      // SESSION_ID = crypto.randomUUID() genere une fois au chargement.
      // Il permet au serveur de maintenir l'historique de conversation.
      body: JSON.stringify({ query: q, key: API_KEY, session_id: SESSION_ID }),
      signal: abortCtrl.signal  // Lie ce fetch a l'AbortController
    });

    typing.remove();
    const bot = addBubble("bot");  // Bulle de reponse, initialement vide
    const reader = res.body.getReader();
    const decoder = new TextDecoder();
    let buf = "", fullText = "";

    // Boucle de lecture du stream SSE ligne par ligne
    while (true) {
      const { done, value } = await reader.read();
      if (done) break;
      // Decode le chunk binaire en texte (stream: true = mode partiel)
      buf += decoder.decode(value, { stream: true });
      const lines = buf.split("\n"); buf = lines.pop();

      for (const line of lines) {
        if (!line.startsWith("data:")) continue;  // Ignore les lignes vides
        const data = JSON.parse(line.slice(5).trim());

        if (data.token) {
          // Ajout progressif du token a la bulle de reponse
          fullText += data.token;
          bot.textContent = fullText;
          chat.scrollTop = chat.scrollHeight;  // Auto-scroll vers le bas
        } else if (data.error) {
          bot.className = "arb-bubble warn";
          bot.textContent = "Je n'ai pas trouve d'information.";
        } else if (data.done) {
          // Fin du stream : affichage des sources comme puces cliquables
          addSources(data.sources || []);
        }
      }
    }
  } catch (err) {
    typing.remove();
    // AbortError = arret volontaire par l'utilisateur : pas d'erreur a afficher.
    // Toute autre erreur (reseau, serveur) est affichee en bulle d'avertissement.
    if (err.name !== "AbortError") {
      addBubble("warn").textContent = "Erreur de connexion.";
    }
  }
  abortCtrl = null;
  setStopMode(false);  // Retour au bouton Envoyer normal
  inputEl.focus();
}
\end{lstlisting}

\textbf{Explication :} L'utilisation de \texttt{AbortController} est une d\'{e}cision technique
importante : elle permet \`{a} l'utilisateur d'interrompre une g\'{e}n\'{e}ration longue en cliquant
sur le bouton Stop (carr\'{e} rouge). L'annulation du \texttt{fetch} envoie un signal
\texttt{AbortError} que le bloc \texttt{catch} intercepte sans afficher d'erreur.
Le texte d\'{e}j\`{a} affich\'{e} est conserv\'{e}, ce qui est plus agr\'{e}able que d'effacer la
r\'{e}ponse partielle. La lecture ligne par ligne avec un buffer (\texttt{buf}) est n\'{e}cessaire
car les chunks TCP re\c{c}us ne correspondent pas forc\'{e}ment aux \'{e}v\'{e}nements SSE :
plusieurs tokens peuvent arriver dans un m\^{e}me chunk, ou un \'{e}v\'{e}nement peut \^{e}tre
coup\'{e} en deux chunks successifs.

% ============================================================
\chapter{Synth\`{e}se technique -- D\'{e}cisions et justifications}

\begin{center}
\renewcommand{\arraystretch}{1.5}
\begin{tabularx}{\linewidth}{>{\bfseries}p{3.5cm} X p{2.5cm}}
  \toprule
  \rowcolor{rowblue}
  \textbf{D\'{e}cision} & \textbf{Justification} & \textbf{Fichier} \\
  \midrule
  ChromaDB local & Seule BD vectorielle satisfaisant co\^{u}t z\'{e}ro + hors ligne + confidentialit\'{e} & config.py \\
  Ollama + Mistral & Stack 100\% gratuit, offline, optimis\'{e} fran\c{c}ais & config.py \\
  Docling en mode par d\'{e}faut & Meilleure pr\'{e}cision sur les brochures touristiques complexes & config.py \\
  MarkdownChunker pour Docling & Pr\'{e}serve les entit\'{e}s s\'{e}mantiques (1 h\^{o}tel = 1 chunk) & pipeline.py \\
  Cascade de fallbacks PDF & Robustesse face \`{a} la diversit\'{e} des documents & loaders.py \\
  Recherche hybride dense+keyword & Combine pr\'{e}cision s\'{e}mantique et correspondance exacte des noms & indexing.py \\
  temperature=0.0 & R\'{e}ponses d\'{e}terministes, essentielles pour un assistant factuel & generation.py \\
  Streaming SSE & Exp\'{e}rience utilisateur progressive, pas d'attente & server.py \\
  AbortController JS & Interruptibilit\'{e} de la g\'{e}n\'{e}ration par l'utilisateur & widget.js \\
  IIFE widget & Auto-contenu, int\'{e}grable sur tout site sans d\'{e}pendances & widget.js \\
  \bottomrule
\end{tabularx}
\end{center}

\section{Points d'am\'{e}lioration identifi\'{e}s}

\begin{enumerate}[leftmargin=*]
  \item \textbf{Passage \`{a} Llama 3 ou Phi-3} : des tests pr\'{e}liminaires sugg\`{e}rent
        que Llama 3 (8B) produit des r\'{e}ponses plus d\'{e}taill\'{e}es en fran\c{c}ais.
        Une \'{e}valuation comparative via Ragas est planifi\'{e}e en S16.
  \item \textbf{Cache de r\'{e}ponses} : les questions fr\'{e}quentes (horaires, tarifs standards)
        pourraient \^{e}tre mises en cache pour r\'{e}duire la latence et la charge sur Ollama.
  \item \textbf{Persistance des sessions} : les sessions en m\'{e}moire sont perdues au red\'{e}marrage.
        Un stockage SQLite l\'{e}ger am\'{e}liorerait la continuit\'{e} pour les sessions longues.
\end{enumerate}

% ============================================================
\chapter{Bilan hebdomadaire -- Semaines 11 \`{a} 15}

\section{Semaine 11 (29 avril -- 5 mai) -- Initialisation}
\begin{infobox}
\textbf{Th\`{e}me} : Analyse des besoins, choix des technologies, mise en place de l'environnement
de d\'{e}veloppement (Ollama, ChromaDB, Flask, Docling).
\end{infobox}
\begin{itemize}[leftmargin=*]
  \item R\'{e}union avec le ma\^{i}tre de stage pour d\'{e}finir le p\'{e}rim\`{e}tre du projet.
  \item Installation et test des technologies candidates (Ollama, ChromaDB, PDFPlumber, Docling).
  \item Premier prototype de pipeline d'ingestion avec un PDF touristique test.
  \item Mise en place de la structure du projet (\texttt{legal\_rag/}, \texttt{server.py}).
\end{itemize}

\section{Semaine 12 (6 -- 12 mai) -- Loaders et chunking}
\begin{infobox}
\textbf{Th\`{e}me} : D\'{e}veloppement des 7 modes d'extraction PDF, impl\'{e}mentation de
StructuralChunker et MarkdownChunker.
\end{infobox}
\begin{itemize}[leftmargin=*]
  \item Impl\'{e}mentation de PDFLoader avec support des modes : docling, markdown, pdfplumber,
        hybrid, ocr, vision, pymupdf.
  \item D\'{e}veloppement de StructuralChunker (d\'{e}coupage par paragraphes) et MarkdownChunker
        (d\'{e}coupage sur les titres \#\#).
  \item Probl\`{e}me identifi\'{e} : les prix flottants dans les brochures Docling se retrouvaient
        d\'{e}tach\'{e}s de leur h\^{o}tel. Solution : algorithme de redistribution par heuristique.
  \item Tests d'ingestion sur le corpus tourisme Centre-Corse (brochures h\'{e}bergement).
\end{itemize}

\section{Semaine 13 (13 -- 19 mai) -- Retrieval et scraping web}
\begin{infobox}
\textbf{Th\`{e}me} : Pipeline de retrieval hybride (dense + keyword), scraping du site
\texttt{tourisme-centrecorse.corsica}, mode WEB\_MODE separate/mixed.
\end{infobox}
\begin{itemize}[leftmargin=*]
  \item Impl\'{e}mentation de la recherche hybride : embedding dense + keyword sur entit\'{e}s nomm\'{e}es.
  \item D\'{e}veloppement de WebLoader avec parsing sitemap.xml et extraction schema.org.
  \item Impl\'{e}mentation des modes WEB\_MODE (\texttt{separate} : docs en priorit\'{e} avec fallback web ;
        \texttt{mixed} : fusion des deux sources).
  \item Analyse du comportement des deux modes sur des questions test (comparatif doc\'{e}s).
\end{itemize}

\section{Semaine 14 (20 -- 26 mai) -- Serveur et frontend}
\begin{infobox}
\textbf{Th\`{e}me} : Serveur Flask avec streaming SSE, widget JS int\'{e}grable, interface
standalone avec design Figma.
\end{infobox}
\begin{itemize}[leftmargin=*]
  \item Impl\'{e}mentation de la route \texttt{/ask} avec Server-Sent Events et gestion des sessions.
  \item D\'{e}veloppement du widget JavaScript int\'{e}grable (\texttt{widget.js}) avec
        FAB, streaming SSE, AbortController, et bouton stop.
  \item Impl\'{e}mentation de l'interface standalone (\texttt{index.html}) selon la maquette
        Figma : widget flottant, mode plein \'{e}cran, responsive mobile.
  \item Int\'{e}gration du syst\`{e}me multi-client avec clés API distinctes (\texttt{clients.py}).
\end{itemize}

\section{Semaine 15 (27 mai -- 5 juin) -- Tests et \'{e}valuation}
\begin{infobox}
\textbf{Th\`{e}me} : Tests fonctionnels, \'{e}valuation qualitative avec tableau Word, r\'{e}glage
des param\`{e}tres RAG (MAX\_DISTANCE, n\_results, WEB\_MODE).
\end{infobox}
\begin{itemize}[leftmargin=*]
  \item Cr\'{e}ation du tableau d'\'{e}valuation Word (\texttt{evaluation\_rag.docx}) pour noter
        les r\'{e}sultats selon diff\'{e}rentes configurations de param\`{e}tres.
  \item Tests sur 8 configurations : variations de \texttt{n\_results} (4, 8, 12),
        \texttt{max\_distance} (0.6, 0.9, 1.2), mod\`{e}le (mistral, llama3), retriever.
  \item Analyse des divergences entre les modes \texttt{separate} et \texttt{mixed}
        (certaines questions obtiennent des r\'{e}ponses dans un mode mais pas l'autre).
  \item Configuration finale retenue : \texttt{n\_results=12}, \texttt{max\_distance=0.9},
        \texttt{PDF\_MODE=docling}, \texttt{WEB\_MODE=mixed}.
\end{itemize}

% ==================== APPENDIX ====================
\appendix
\chapter*{Annexes}
\addcontentsline{toc}{chapter}{Annexes}

\section*{A. Stack technologique complet}
\begin{center}
\renewcommand{\arraystretch}{1.4}
\begin{tabularx}{\linewidth}{>{\ttfamily}p{3.5cm} p{2cm} X}
  \toprule
  \rowcolor{rowblue}
  \textbf{Package} & \textbf{Version} & \textbf{R\^{o}le dans le projet} \\
  \midrule
  flask        & 3.x   & Serveur REST, routes API, streaming SSE \\
  flask-cors   & 4.x   & Autorisation des appels depuis d'autres domaines \\
  chromadb     & 0.5.x & Base de donn\'{e}es vectorielle locale persistante \\
  ollama       & 0.2.x & Client Python pour Ollama (embeddings + g\'{e}n\'{e}ration) \\
  pymupdf      & 1.24.x& Extraction texte PDF (mode pymupdf et hybrid) \\
  pymupdf4llm  & 0.0.x & Conversion PDF vers Markdown structur\'{e} \\
  docling      & 2.x   & Extraction layout PDF par deep learning (IBM) \\
  pdfplumber   & 0.11.x& Extraction de tableaux et colonnes PDF \\
  pytesseract  & 0.3.x & OCR Tesseract pour les PDFs scann\'{e}s \\
  beautifulsoup4 & 4.x & Extraction de contenu web (scraping) \\
  requests     & 2.x   & Appels HTTP pour le scraping web \\
  ragas        & 0.1.x & \'{E}valuation automatique du pipeline RAG \\
  \bottomrule
\end{tabularx}
\end{center}

\section*{B. R\'{e}f\'{e}rences}
\begin{itemize}[leftmargin=*]
  \item Lewis et al. (2020). \textit{Retrieval-Augmented Generation for Knowledge-Intensive NLP Tasks}.
        \url{https://arxiv.org/abs/2005.11401}
  \item ChromaDB -- Documentation officielle : \url{https://docs.trychroma.com}
  \item Ollama -- Documentation officielle : \url{https://ollama.com/docs}
  \item IBM Docling : \url{https://github.com/DS4SD/docling}
  \item Ragas -- \'{E}valuation RAG : \url{https://docs.ragas.io}
  \item Flask -- Server-Sent Events : \url{https://flask.palletsprojects.com/en/stable/patterns/streaming/}
\end{itemize}

\end{document}
